% © 2026 Ing. David Jaroš — CC BY-NC-ND 4.0
%
% This work is licensed under a Creative Commons Attribution-NonCommercial-NoDerivatives
% 4.0 International License (CC BY-NC-ND 4.0).
%
% License History: Earlier drafts (up to v0.3) were released under CC BY 4.0.
% From v0.4 onward, all material is released under CC BY-NC-ND 4.0 to protect
% the integrity of the theoretical work during ongoing academic development.
%
% See LICENSE.md for full license text.
%
% File: research_tracks/T3_ALPHA/hecke_alpha_connection.tex
%
% Purpose: Document the investigation of the Hecke eigenvalue coincidence
%   a_{137}(N=76, k=2) = -11 = -genus(X_0(137)), and its implications for
%   the alpha derivation programme.
%
% Numerical results generated by: tools/hecke_alpha_scan.py (2026-05-10)
%
% Data sources: reports/gamma0_137_invariants.md §3,
%               reports/prime_137_structural_audit.md §1.3

\ifdefined\INCLUDEMODE
\else
  \documentclass[12pt,a4paper]{article}
  \usepackage[a4paper, margin=2.5cm]{geometry}
  \usepackage{amsmath, amssymb, amsthm}
  \usepackage{mathtools}
  \usepackage{hyperref}
  \usepackage{xcolor}
  \usepackage{booktabs}
  \usepackage{longtable}
  \usepackage{mdframed}
  \usepackage{tcolorbox}
  \tcbuselibrary{skins,breakable}

  \newtheorem{theorem}{Theorem}[section]
  \newtheorem{lemma}[theorem]{Lemma}
  \newtheorem{proposition}[theorem]{Proposition}
  \theoremstyle{definition}
  \newtheorem{definition}[theorem]{Definition}
  \theoremstyle{remark}
  \newtheorem{remark}[theorem]{Remark}

  \newmdenv[backgroundcolor=yellow!15, linecolor=orange!70, linewidth=1.5pt]{gapbox}
  \newmdenv[backgroundcolor=green!12, linecolor=green!60!black, linewidth=1.5pt]{resultbox}
  \newmdenv[backgroundcolor=blue!6, linecolor=blue!45, linewidth=1.5pt]{insightbox}
  \newmdenv[backgroundcolor=red!8, linecolor=red!60, linewidth=1.5pt]{warnbox}

  \newtcolorbox{statusbox}[1][]{colback=gray!8!white,colframe=black!60,
    arc=3pt,boxrule=1pt,fonttitle=\bfseries,title=Status Summary,#1}

  \title{\textbf{Hecke Eigenvalue Connection to $\alpha$}\\[4pt]
         \large Track T3\_ALPHA: Approach B --- $a_{137} = \mathrm{genus}(X_0(137)) = 11$}
  \author{Ing.\ David Jaro\v{s}}
  \date{2026-05-10}

  \begin{document}
  \maketitle
\fi

%──────────────────────────────────────────────────────────────────────────────
\section{Setup and Data}
\label{sec:hecke:setup}
%──────────────────────────────────────────────────────────────────────────────

\subsection{Hecke eigenvalues at $p = 137$}

From \texttt{reports/gamma0\_137\_invariants.md}~\S{}3, the following Hecke
eigenvalues $a_{137}$ of modular forms at coprime levels are known:

\begin{center}
\begin{tabular}{rrrcc}
  \toprule
  Level $N$ & Weight $k$ & $a_{137}$ & Weil bound $2p^{(k-1)/2}$ & Within bound? \\
  \midrule
  76  & 2 & $-11$ & $23.41$ & YES \\
  7   & 4 & $+2274$ & $3207.09$ & YES \\
  208 & 6 & $-38286$ & $439371.03$ & YES \\
  \bottomrule
\end{tabular}
\end{center}

\noindent The form at level $N=76$, weight $k=2$ corresponds (heuristically) to the
electron generation; weight $k=4$ to the muon; weight $k=6$ to the tau.

\subsection{Genus of $X_0(137)$}

From the Riemann--Hurwitz formula for prime level $p$
(Diamond \& Shurman, Theorem~3.1.1):
\begin{align}
  \mu(\Gamma_0(137)) &= 138 ,\\
  \nu_2(137) &= 1 + \left(\tfrac{-4}{137}\right) = 2 , \quad (137 \equiv 1 \bmod 4)\\
  \nu_3(137) &= 1 + \left(\tfrac{-3}{137}\right) = 0 , \quad (137 \equiv 2 \bmod 3)\\
  g(X_0(137)) &= 1 + \frac{138}{12} - \frac{2}{4} - \frac{0}{3} - \frac{2}{2} = 11 .
  \label{eq:genus137}
\end{align}

\noindent\textbf{Observation}:
$|a_{137}(N=76, k=2)| = |-11| = 11 = g(X_0(137))$.

%──────────────────────────────────────────────────────────────────────────────
\section{Test: Is $a_p = \pm\,\mathrm{genus}(X_0(p))$ structural?}
\label{sec:hecke:test}
%──────────────────────────────────────────────────────────────────────────────

\subsection{Atkin-Lehner obstruction at prime level}

For a newform $f \in S_2(\Gamma_0(p))$ at \emph{prime} level $p$, the Atkin--Lehner
involution $w_p$ acts as $w_p f = \varepsilon_p f$ with $\varepsilon_p = \pm 1$.
The Hecke eigenvalue at $p$ satisfies
\begin{equation}
  a_p(f) = \varepsilon_p \in \{-1, +1\} \quad
  \text{(for weight 2, prime level)}.
  \label{eq:atkin_lehner}
\end{equation}
Therefore $|a_p(f_{\mathrm{prime\,level}})| = 1$, \emph{not} $g(X_0(p))$.
The coincidence $|a_{137}| = g = 11$ can only arise for a form $f$ at a
level $N$ \emph{coprime to} $137$ (a ``good prime'').

\subsection{Primes $p$ with $g(X_0(p)) = 11$}

The scan over $p \in [50, 300]$ (from \texttt{tools/hecke\_alpha\_scan.py}) gives:

\begin{center}
\begin{tabular}{rrrr}
  \toprule
  $p$ & $\nu_2$ & $\nu_3$ & $g(X_0(p))$ \\
  \midrule
  131 & 0 & 0 & 11 \\
  \textbf{137} & \textbf{2} & \textbf{0} & \textbf{11} \\
  139 & 0 & 2 & 11 \\
  149 & 2 & 0 & 12 \\
  \bottomrule
\end{tabular}
\end{center}

There are \textbf{three} primes in $[50,300]$ with $g(X_0(p)) = 11$:
$p \in \{131, 137, 139\}$.  The form at $N=76$, $k=2$ has $a_{137} = -11$.
Whether $a_{131}$ or $a_{139}$ (of the same form) also equal $\pm 11$ is an
open question requiring direct computation from the LMFDB database or
Sage/Magma.

\subsection{Twin-prime lookup status (Task 2 update)}

The dedicated lookup note
\texttt{reports/hecke\_eigenvalue\_twin\_prime\_test.md} records the latest
attempt (2026-05-11) to determine $a_{131}(76.2.a.a)$ and $a_{139}(76.2.a.a)$.

Current state:
\begin{itemize}
  \item $a_{137} = -11$ is known and confirmed from previous runs.
  \item $a_{131}$ and $a_{139}$ remain unresolved in this sandbox because:
    (i) LMFDB network access is unavailable (DNS failure), and
    (ii) SageMath is not installed locally.
\end{itemize}

Therefore the binary test
``$|a_{131}|=11$ or $|a_{139}|=11$?'' is still \textbf{[OPEN]} and the
classification of the $|a_{137}|=g=11$ signal remains provisional.

\subsection{Structural vs.\ coincidental assessment}

\begin{insightbox}
\textbf{No theorem} of the form ``$a_p(f_{N=76}) = \pm g(X_0(p))$ for all
primes $p$ with $\gcd(N, p) = 1$'' is known in the literature.  Such a
relation would be extraordinary: it would equate the Hecke eigenvalue of a
fixed form at a varying prime $p$ with a topological invariant of a different
modular curve.

The coincidence $|a_{137}(N=76, k=2)| = g(X_0(137)) = 11$ is therefore
classified as \textbf{[NC]} (numerical coincidence) or at best \textbf{[MC]}
(motivated conjecture) pending computation of $a_{131}$ and $a_{139}$ for
the same form.
\end{insightbox}

%──────────────────────────────────────────────────────────────────────────────
\section{Lepton Mass Ratios from Hecke Eigenvalues}
\label{sec:hecke:lepton}
%──────────────────────────────────────────────────────────────────────────────

The lepton mass ratio reconstruction uses the ratios of $|a_{137}|$ across
weights:
\begin{align}
  R_\mu &= \frac{|a_{137}(k=4)|}{|a_{137}(k=2)|}
         = \frac{2274}{11} = 206.727 , \quad
         \text{exp: } 206.768 , \quad \text{err: } 0.020\% , \\
  R_\tau &= \frac{|a_{137}(k=6)|}{|a_{137}(k=2)|}
          = \frac{38286}{11} = 3480.55 , \quad
          \text{exp: } 3477.23 , \quad \text{err: } 0.095\% .
\end{align}

This signal is:
\begin{enumerate}
  \item \textbf{Independent of $V_{\mathrm{eff}}$}: uses no winding potential.
  \item \textbf{Prime-specific}: $p=137$ is the unique prime in $[50,300]$ achieving
        both errors $< 0.1\%$ (established by scan, see \texttt{reports/prime\_137\_structural\_audit.md}~\S{}1.3).
  \item \textbf{Weil-bounded}: all three eigenvalues satisfy $|a_{137}| \leq 2p^{(k-1)/2}$.
\end{enumerate}

However, the forms at levels $76, 7, 208$ were found by search; their derivation
from UBT first principles (why these levels and weights?) remains an open problem.

%──────────────────────────────────────────────────────────────────────────────
\section{Hecke Structure and $\alpha^{-1} = 137$}
\label{sec:hecke:alpha}
%──────────────────────────────────────────────────────────────────────────────

\subsection{The proposed connection}

The conjecture is that $\alpha^{-1} = p$ is fixed by the Hecke structure
through the requirement that the three ``generational'' modular forms
(at levels $N=76, 7, 208$ and weights $k=2, 4, 6$) simultaneously satisfy
the lepton mass ratio conditions at prime $p$.  The prime $p$ where this
``resonance'' is global minimum is $p=137$.

\subsection{Test: Is $a_p/p = 11/137$ unique?}

If $|a_p(N=76, k=2)| = 11$ for all three primes $p \in \{131, 137, 139\}$
with $g(X_0(p))=11$, then the ratio $a_p/p$ differs:

\begin{center}
\begin{tabular}{rrrr}
  \toprule
  $p$ & $g(p)$ & $11/p$ (hypothetical if $|a_p|=11$) & vs.\ $11/137$ \\
  \midrule
  131 & 11 & $0.08397$ & $+0.00368$ \\
  \textbf{137} & \textbf{11} & \textbf{0.08029} & \textbf{reference} \\
  139 & 11 & $0.07914$ & $-0.00116$ \\
  \bottomrule
\end{tabular}
\end{center}

\noindent Whether $a_p$ actually equals $11$ for $p=131$ and $p=139$ is not
computationally accessible here.

\subsection{Verdict on the Hecke $\to \alpha$ connection}

\begin{warnbox}
\textbf{Status}: The connection ``Hecke eigenvalue $\to \alpha^{-1}$'' is
\textbf{[OPEN]}.

The lepton mass ratio signal at $p=137$ is a strong numerical result [MC]
and is prime-specific (global minimum in $[50,300]$).  However:
\begin{itemize}
  \item The forms at $N=76, 7, 208$ were found by search, not derived from UBT.
  \item The coincidence $|a_{137}|=g=11$ is [NC/MC], not [L1].
  \item The mechanism connecting Hecke eigenvalues to $\alpha^{-1}$ via
        the path integral is currently at no-go status (see
        \texttt{reports/hecke\_path\_integral\_no\_go\_or\_success.md}).
\end{itemize}
\end{warnbox}

%──────────────────────────────────────────────────────────────────────────────
\begin{statusbox}
\textbf{$a_{137} = \mathrm{genus}(X_0(137))$?}: ANO numericky ($|a_{137}|=11=g$).
Ale není to strukturální fakt odvozitelný z Atkin-Lehnerovy teorie.

\textbf{Je shoda unikátní pro $p=137$?}: NEZNÁMÉ — vyžaduje LMFDB data
pro $p=131$ a $p=139$ (stejná forma $N=76$, $k=2$).  Pokud $a_{131}=a_{139}=-11$
také, pak shoda není unikátní.

\textbf{Spojení Hecke $\to \alpha$}: OPEN — lepton-mass signál je silný [MC];
Hecke path-integral mechanismus je na no-go statusu.

\textbf{Klasifikace}: [NC/MC] pro shodu $|a_{137}|=g=11$; [MC] pro lepton mass signal;
[OPEN] pro celkové spojení Hecke $\to \alpha^{-1}$.
\end{statusbox}

\ifdefined\INCLUDEMODE
\else
  \end{document}
\fi
